% Copyright 2026 Yilin Liu and GPT-6 Astra Pro. Apache-2.0.
\documentclass[11pt]{article}
\usepackage[a4paper,margin=26mm]{geometry}
\usepackage[T1]{fontenc}
\usepackage{amsmath,amssymb,amsthm,mathtools}
\usepackage[numbers]{natbib}
\usepackage[colorlinks=true]{hyperref}
\newtheorem{theorem}{Theorem}[section]
\newtheorem{lemma}[theorem]{Lemma}
\newtheorem{corollary}[theorem]{Corollary}
\newtheorem{proposition}[theorem]{Proposition}
\newcommand{\poly}{\operatorname{poly}}
\newcommand{\polylog}{\operatorname{polylog}}
\title{Integral Character Profiles and Ordered Crossing Circuits}
\author{Yilin Liu and GPT-6 Astra Pro}
\date{28 October 2026}
\begin{document}
\maketitle
The selected result combines a normalization-one integral crossing,
simultaneous ordered insertion and shared matrix products with a full
three-shape integer matrix bound. The recent integral spectrum,
convolution, separation and sector theorems retain their explicit source
attribution. The weighted field-profile analysis is also prior work;
the composition here preserves one integer scheme chosen before every
associative unital coefficient ring. This manuscript retains the completed
mathematical argument from the supplied checkpoint, without claiming
external peer review, historical priority or a complete local formal
verification. No satisfiability algorithm is a premise. We give one
selected $O(N^{1.719516})$ arithmetic circuit bound rather than a list
of numerically weaker substitutions or direct SAT applications.
\tableofcontents
\section{Finite crossing and the integral lattice}
% Copyright 2026 Yilin Liu and GPT-6 Astra Pro. Apache-2.0.
% Curated from sat_research_checkpoint.zip: paper/elementary_crossing_operator_lemma.tex
% Complete finite proof retained; no SAT dependency or formal receipts imported.
\subsection{An elementary proof of the crossing operator identity}
\label{sec:elementary-crossing-operator}

We give a self-contained proof of the finite operator identity used in
the integral crossing construction. This isolates the algebraic
consequence from the satisfiability algorithm: no claim about that
algorithm is needed in this subsection. The identity is the one in
Appendix~B, Lemma~9 of~\cite{anthropic2026}; the proof below uses only
normal ordering and a finite matching expansion.

Let $V=\mathbb Q[p_1,p_2,\ldots]$, graded by $\deg p_j=j$, and let
$V_E$ denote the homogeneous component of weight $E$; set $V_E=0$
for $E<0$. For $k\ge1$ define
\[
 D_{-k}f=p_kf,\qquad D_kf=k\frac{\partial f}{\partial p_k}.
\]
Thus $[D_k,D_{-l}]=k[k=l]\operatorname{Id}$, while creation
operators commute with one another, as do annihilation operators.
For $m\ge1$ write
\begin{align*}
 J_m^+(z)&=\sum_{l\ge1}\binom{l-1}{m-1}D_{-l}z^{l-m},\\
 J_m^-(z)&=\sum_{k\ge1}\binom{-k-1}{m-1}D_kz^{-k-m}.
\end{align*}
The first sum actually starts at $l=m$. For a partition
$\mu=(\mu_1,\ldots,\mu_t)$, define
\[
 Y(p_\mu,z)=:\!\prod_{i=1}^t
       (J_{\mu_i}^+(z)+J_{\mu_i}^-(z))\!:
       =\sum_{r\in\mathbb Z}p_\mu[r]z^{r-|\mu|}.
\]
The colons mean that in each term all creation operators are placed
to the left of all annihilation operators, without adding commutator
terms. Extend the modes linearly to $V$. Every $a[r]$ maps
$V_E$ to $V_{E+r}$. Its action on a polynomial is well-defined:
the positive oscillator indices in a nonzero term have sum at most
the input weight, and the prescribed mode then bounds the negative
indices as well.

\begin{lemma}[Finite crossing operator identity]
\label{lem:elementary-crossing-operator}
Let $h\ge1$, $b\in V_h$, and $a=\sum_{h'=0}^h a_{h'}$ with
$a_{h'}\in V_{h'}$. Set
\[
 a^{\langle i\rangle}=\sum_{h'=0}^h\binom{h+h'}i a_{h'},\qquad
 Q_s(a,b)=\sum_{i=0}^{2h}(-1)^i
    a^{\langle i\rangle}[s-h+i]b[h-s-i].
\]
On every $V_E$ the following statements hold:
\begin{enumerate}
\item Every product in the definition of $Q_s$ can be reversed
simultaneously without changing $Q_s$.
\item $Q_s=0$ if $|s|>E+h$.
\item $\sum_{|s|\le E+h}Q_s=(a[-h]b)\operatorname{Id}_{V_E}$.
\end{enumerate}
\end{lemma}

\begin{proof}
Linearity reduces the claim to $a=p_\mu$ of weight $h'\le h$
and $b=p_\nu$ of weight $h$. Put $H=h+h'$. Let $\iota_+$
mean expansion in nonnegative powers of $w/z$, and $\iota_-$
expansion in nonnegative powers of $z/w$. Directly from the
commutation relations,
\begin{equation}
\label{eq:elementary-contraction}
 [J_m^-(z),J_n^+(w)]
 =\iota_+\frac{\gamma_{m,n}}{(z-w)^{m+n}},\qquad
 \gamma_{m,n}=(-1)^{m-1}
       \frac{(m+n-1)!}{(m-1)!(n-1)!}.
\end{equation}
Indeed, for $m=n=1$ the contraction is
$\sum_{k\ge1}kz^{-k-1}w^{k-1}=\iota_+(z-w)^{-2}$.
Applying the divided derivatives
$\partial_z^{m-1}\partial_w^{n-1}/((m-1)!(n-1)!)$
gives~\eqref{eq:elementary-contraction}. Reversing the fields gives
the expansion $\iota_-$ of the same rational function, since
$(w-z)^{-2}=(z-w)^{-2}$ before differentiation.

Label equal parts of $\mu$ and $\nu$ separately. Commuting
annihilation operators of the first field past creation operators
of the second gives a sum over partial matchings between these two
lists of parts. A matched pair $(m,n)$ contributes the scalar
in~\eqref{eq:elementary-contraction}; all unmatched operators form
one completely normally ordered product. This assertion follows by
repeatedly using $AB=BA+[A,B]$: every commutator is scalar, so it
creates no further noncentral terms, and each labelled operator can
be contracted at most once.

Fix an input and an output monomial of weight $E$. The matrix
coefficient of any completely normally ordered product of the
unmatched operators is a Laurent polynomial in $z,w$. Its
annihilation indices sum to at most $E$. Since input and output
weights agree, its creation indices have that same sum. Hence all
indices are bounded by $E$, and only finitely many terms occur.
The original two ordered products have well-defined matrix
coefficients in the corresponding iterated Laurent-series rings:
the rightmost field has powers bounded below, and fixing one such
power bounds its intermediate weight. The preceding commutations
may therefore also be justified coefficientwise by truncating
oscillator indices and then taking stabilized coefficients.

For a partial matching let $D$ be the sum of the weights of all
matched parts, so $D\le H$. Multiplying its contraction factors
by $(z-w)^H$ leaves $(z-w)^{H-D}$ times an ordinary Laurent
polynomial. Consequently
\[
 (z-w)^H [u]Y(a,z)Y(b,w)v
 \quad\hbox{and}\quad
 (z-w)^H [u]Y(b,w)Y(a,z)v
\]
are the two expansions of one and the same Laurent polynomial
$F_{u,v}(z,w)$. Here $u,v$ are the chosen weight-$E$ monomials,
and $[u]$ extracts a coefficient. This polynomial is homogeneous
of total degree zero, hence has the form
$\sum_s F_s z^s w^{-s}$ with finite support.

Reading the coefficient $z^sw^{-s}$ in the first order gives
\[
 F_s=\sum_{i=0}^H(-1)^i\binom Hi
       [u]a[s-h+i]b[h-s-i]v=[u]Q_s(a,b)v.
\]
Reading it in the other order reverses every product. This proves
the first assertion. In the first order the rightmost mode is
$b[-r]$ with $r=s-h+i$, so it vanishes on $V_E$ when $r>E$.
Thus every term vanishes for $s>E+h$. In the reverse order the
rightmost mode is $a[r]$, which vanishes when $r<-E$. Since
$r\le s+h'\le s+h$, every term vanishes for $s<-E-h$.
This proves the second assertion.

At $z=w=1$, every partial matching with an unmatched part vanishes,
because then $H-D>0$. Fully matched terms leave only a scalar
times the identity. Denote the sum of their scalar coefficients by
$c(a,b)$. It is independent of $E$, so
\[
 \sum_s Q_s(a,b)=c(a,b)\operatorname{Id}_{V_E}.
\]
To identify the scalar, use the vacuum vector $1\in V_0$.
Its two-field matrix coefficient contains only fully matched terms,
and is therefore $\iota_+c(a,b)(z-w)^{-H}$. On the other hand,
$Y(b,w)1=b+O(w)$: each creation summand starts with its own
$p_n$ at power zero. Taking the coefficient of $w^0$ and then
the constant polynomial coefficient gives
$a[-h]b\,z^{-H}=c(a,b)z^{-H}$. Hence $c(a,b)=a[-h]b$,
which proves the third assertion. Summing homogeneous components
of $a$ gives the stated generality.
\end{proof}

\begin{lemma}[The finite binomial duality]
\label{lem:elementary-crossing-duality}
Define the multiplicative linear transform $T:V\to V$ by
\[
 T(p_n)=(-1)^{n-1}\sum_{j=1}^n\binom njp_j.
\]
For partitions $\mu,\nu$ of $h$, let
$z_\mu=\prod_{k\ge1}k^{m_\mu(k)}m_\mu(k)!$. Then
\[
 T(p_\mu)[-h]p_\nu=z_\mu[\mu=\nu].
\]
\end{lemma}

\begin{proof}
This is the unscaled duality in the source's Lemma~10, which we
rederive to make the dependency explicit. Vandermonde's identity,
including negative upper binomial arguments, gives
\[
 \sum_{j=1}^n\binom nj\binom{-k-1}{j-1}
   =\binom{n-k-1}{n-1}.
\]
It is the coefficient of degree $n-1$ in
$(1+t)^n(1+t)^{-k-1}$. The result is zero for $1\le k<n$
and equals $(-1)^{n-1}$ for $k=n$. The sign in $T(p_n)$
therefore makes its surviving coefficient at $k=n$ equal to one.

In a mode of total shift $-h$ acting on $V_h$, any creation
operator would require total differentiation weight greater than
$h$, so it vanishes. In the remaining all-annihilation terms of
$T(p_\mu)[-h]$, each index $k_i$ is at least $\mu_i$, and their
sum is $h=\sum_i\mu_i$. Hence every $k_i=\mu_i$, and
\[
 T(p_\mu)[-h]p_\nu
   =\left(\prod_i D_{\mu_i}\right)p_\nu
   =z_\mu[\mu=\nu].
\]
If all derivative multiplicities fit those of $p_\nu$, equality
of the total weights forces the two partitions to be identical;
otherwise a derivative annihilates the monomial.
\end{proof}

The argument is first carried out over $\mathbb Q$ to define the
displayed operators without changing bases. When an integral lattice
is preserved by their modes, the resulting equality of integral
matrices specializes to every commutative coefficient ring. The
Schur-lattice construction below provides exactly
that lattice and a dual family with scalar $a[-h]b$ equal to a
Kronecker delta.

% Copyright 2026 icekylinx and GPT-6 Astra Pro. Apache-2.0.
% Curated from sat_research_checkpoint.zip: paper/integral_schur_extension.tex
% Complete integral crossing theorem and proof retained; unrelated downstream corollaries and duplicate streaming implementation omitted.
\subsection{An integral Schur basis removes the factorial normalization}
\label{sec:integral-schur}

The factorial $h!$ in the source's monomial-basis construction is
not intrinsic to its rank bound. A change to an integral Schur
basis removes it. Integral forms of Heisenberg and lattice vertex
operators are established theory~\cite{huangjing2022}; the proof
below gives the specific finite coefficient statement needed here.
The required source operator identities are independently proved
in Lemmas~\ref{lem:elementary-crossing-operator}
and~\ref{lem:elementary-crossing-duality}. This algebraic result
does not depend on the source SAT algorithm.

\begin{theorem}[Integral crossing pieces]
\label{thm:integral-crossing}
For the modes defined in Section~\ref{sec:elementary-crossing-operator}
and all integers $h,E\ge1$, there are
integer crossing pieces indexed by $-E-h\le s\le E+h$ satisfying
\[
 \sum_s A_s(x,y,u,v)=[x=v][y=u],
\]
where $x,v$ range over the partitions of $h$ and $y,u$ over those
of $E$. Both factorizations of every piece have integer entries,
and their ranks are at most
\[
 R_s^{\rm row}=\sum_{i=0}^{2h}p(E+s-h+i),\qquad
 R_s^{\rm col}=\sum_{i=0}^{2h}p(E+h-s-i).
\]
Here $p(j)=0$ for $j<0$. Thus the normalizing scalar is $1$, and
the identities and rank bounds specialize to every commutative
ring. The integers are explicitly constructible with bit lengths
polynomial in $h+E$.
\end{theorem}

\begin{proof}
Identify the source variables $t_j$ with power-sum symmetric
functions $p_j$. Put
\[
 H(u)=\sum_{n\ge0}h_nu^n
     =\exp\!\left(\sum_{n\ge1}p_nu^n/n\right),
 \qquad E(u)=H(-u)^{-1}.
\]
The integral lattice is the symmetric-function ring
$\Lambda_{\mathbb Z}=\mathbb Z[h_1,h_2,\ldots]$. Products
$h_\lambda$ and Schur functions $s_\lambda$ are integral bases.
The Hall pairing has
\[
 \langle p_\lambda,p_\mu\rangle
   =z_\lambda[\lambda=\mu],\qquad
 \langle s_\lambda,s_\mu\rangle=[\lambda=\mu].
\]

First, every source mode of an element of $\Lambda_{\mathbb Z}$
preserves $\Lambda_{\mathbb Z}$. For a homogeneous state $a$ of
weight $w$, write $Y(a,z)=\sum_r a[r]z^{r-w}$ and extend the
state variable linearly. The source's normal ordering places all
multiplications before all differentiations. Direct summation of
its binomial coefficients gives
\[
 Y(H(u),z)f
 =\frac{H(z+u)}{H(z)}\,
 f\left[p_k\mapsto p_k+z^{-k}-(z+u)^{-k}\right].
\]
Indeed the creation coefficient of $p_l$ is
\[
 \sum_{n\ge1}\frac{u^n}{n}\binom{l-1}{n-1}z^{l-n}
 =\frac{(z+u)^l-z^l}{l},
\]
and the coefficient of $\partial_{p_k}$ in the annihilation
exponential is
\[
 k\sum_{n\ge1}\frac{u^n}{n}\binom{-k-1}{n-1}z^{-k-n}
 =z^{-k}-(z+u)^{-k}.
\]
More generally, normal ordering once across all factors yields
\begin{equation}
\label{eq:integral-mode-generating}
 Y\!\left(\prod_{j=1}^qH(u_j),z\right)f
 =\frac{\prod_{j=1}^qH(z+u_j)}{H(z)^q}\,
 f\left[p_k\mapsto p_k+qz^{-k}
                       -\sum_{j=1}^q(z+u_j)^{-k}\right].
\end{equation}
The substitution adds $q$ copies of the letter $z^{-1}$ and
subtracts the letters $(z+u_j)^{-1}$ in the symmetric-function
alphabet. It is integral: on the complete-function generating
series it multiplies by
\[
 \frac{\prod_j(1-(z+u_j)^{-1}t)}{(1-z^{-1}t)^q}.
\]
The prefactor in~\eqref{eq:integral-mode-generating} is integral
because $H(z)^{-1}=E(-z)$. Fixing the multidegree in the $u_j$ and
the weight of $f$ gives a lower Laurent bound in $z$; every
requested coefficient is a finite sum. Extracting a monomial in
the $u_j$ therefore proves mode integrality for $h_\lambda$, and
hence for every integral symmetric function. In particular, all
mode matrices in the Schur basis are integer matrices.

Next define the source's triangular transform
\[
 T(p_n)=(-1)^{n-1}\sum_{j=1}^n\binom njp_j.
\]
It obeys
\[
 T(H(u))=E\!\left(\frac{u}{1+u}\right),\qquad
 T(h_n)=\sum_{j=1}^n(-1)^{n-j}\binom{n-1}{j-1}e_j.
\]
Thus $T$ preserves the integral weight filtration, including
integrality of each homogeneous component. The finite Vandermonde
calculation in Lemma~\ref{lem:elementary-crossing-duality} gives
\[
 (T p_\lambda)[-h]p_\mu=z_\lambda[\lambda=\mu]
 \quad (|\lambda|=|\mu|=h).
\]
Extending over $\mathbb Q$ identifies this pairing with the Hall
pairing. Consequently
\[
 a_\lambda=T(s_\lambda),\quad b_\nu=s_\nu
 \quad\Longrightarrow\quad
 a_\lambda[-h]b_\nu=[\lambda=\nu].
\]

Use these $a_\lambda,b_\nu$ in the source's $Q_s$ and take matrix
entries in the Schur basis of $V_E$. The state
\[
 a_\lambda^{\langle i\rangle}
 =\sum_{w=0}^h\binom{h+w}{i}(a_\lambda)_w
\]
is integral componentwise. The mode-integrality argument therefore
makes every entry of both factorizations an integer.
Lemma~\ref{lem:elementary-crossing-operator} gives
the stated crossing identity with scalar $1$. The two intermediate
weights are $E+s-h+i$ and $E+h-s-i$, so the dimensions give exactly
the asserted ranks. Standard finite symmetric-function basis
conversion, or the generating identities above, constructs all
entries; factorial denominators in an intermediate power-sum
calculation cancel in the final integer coordinates, and all
intermediate weights are $O(h+E)$, giving polynomial bit height.
\end{proof}

% Copyright 2026 icekylinx and GPT-6 Astra Pro. Apache-2.0.
% Curated from sat_research_checkpoint.zip: paper/algebraic_streaming_coefficients.tex
% Complete finite proof retained; no SAT dependency or formal receipts imported.
\subsection{Computing the integral crossing coefficients without large tables}
\label{sec:algebraic-streaming-coefficients}

The following finite algorithm supplies the construction cost needed
by the algebraic circuit consequences.  It uses only the modes and
the integral Schur family already defined above; it has no
tensor-network or satisfiability-algorithm premise.

\begin{lemma}[Explicit computation of integral crossing coefficients]
\label{lem:algebraic-coefficient-computation}
For the weight parameters $h,E\ge1$ of
Theorem~\ref{thm:integral-crossing}, any entry of a crossing piece
or of either of its displayed factorizations can be computed in
$2^{O(\sqrt{h+E})}$ bit operations and
$\operatorname{poly}(h+E)$ space.  All rational intermediates and
the final integer have bit length polynomial in $h+E$.
The implied constants are absolute.
\end{lemma}

\begin{proof}
Every state, input, output, or intermediate partition in the
crossing has weight $O(h+E)$: in particular, the intermediate
weights $E+s-h+i$ and $E+h-s-i$ are at most $2E+2h$ when
$|s|\le E+h$ and $0\le i\le2h$.  Put $W=4(h+E+1)$ to bound
all these weights.  Partitions of weight at most $W$ can be
enumerated in polynomial space, and their number is
$2^{O(\sqrt W)}$.  The latter bound follows directly from Euler's
product evaluated at $\exp(-1/\sqrt W)$.

\paragraph{A power-sum mode entry.}
Write $m_\sigma(j)$ for the multiplicity of part $j$ in a
partition $\sigma$.  Consider first
\[
 [p_\tau]a[r]p_\sigma.
\]
Return zero unless $r=|\tau|-|\sigma|$.  Normal ordering means
that differentiations act before creations.  Enumerate derivative
multiplicities $0\le d_j\le m_\sigma(j)$.  The required creation
multiplicities are forced to be
\[
 c_j=m_\tau(j)-m_\sigma(j)+d_j.
\]
Discard a vector with a negative $c_j$; otherwise
$0\le c_j\le m_\tau(j)$.  The total mode-shift condition is then
automatic.  The differentiation contributes the integer
\[
 D(d)=\prod_{j\ge1}j^{d_j}(m_\sigma(j))_{d_j},
\]
where $(m)_d=m(m-1)\cdots(m-d+1)$.

For a monomial state $a=p_\mu=p_{\mu_1}\cdots p_{\mu_\ell}$,
form the ordinary homogeneous polynomial
\[
 H_\mu(Y)=\prod_{t=1}^{\ell}
       \left(\sum_{k\ne0}\binom{-k-1}{\mu_t-1}Y_k\right).
\]
In this expression retain only positive $k$ in
$\operatorname{supp}(\sigma)$ and negative $k$ with
$-k\in\operatorname{supp}(\tau)$.  No other differentiations can
survive, and no other creations can occur in the requested output.
The contribution of $(d,c)$ is $D(d)$ times the coefficient of
\[
 \prod_jY_j^{d_j}Y_{-j}^{c_j}
\]
in $H_\mu$.  It is zero unless $\sum_j(d_j+c_j)=\ell$.

For an ordinary homogeneous polynomial $H$ of degree $|e|$,
rectangular finite differences extract precisely this coefficient:
\begin{equation}
\label{eq:algebraic-rectangular-differences}
 [Y^e]H
 =\frac1{\prod_k e_k!}
   \sum_{0\le z_k\le e_k}
       (-1)^{|e|-|z|}\prod_k\binom{e_k}{z_k}H(z).
\end{equation}
Indeed every other monomial of degree $|e|$ has an exponent below
its target in at least one coordinate, and is annihilated by that
coordinate's finite difference.  The same identity works with
coefficients in a univariate polynomial ring.

A derivative vector is a subpartition of $\sigma$, so the number
of such vectors is at most
\[
 \prod_j(m_\sigma(j)+1)\le\sum_{v=0}^W p(v)
 =2^{O(\sqrt W)}.
\]
For each vector, the finite-difference grid has at most
\[
 \prod_j(d_j+1)(c_j+1)
 \le\prod_j(m_\sigma(j)+1)(m_\tau(j)+1)
 =2^{O(\sqrt W)}
\]
points.  A fixed product of these bounds remains
$2^{O(\sqrt W)}$.  At a grid point the displayed product of
$\ell\le W$ short linear forms is evaluated directly.

\paragraph{The transformed state and its weight marker.}
For $a=T(p_\mu)$, use instead
\[
 H_\mu(Y,Z)=
 \prod_{t=1}^{\ell}(-1)^{\mu_t-1}
 \left(\sum_{j=1}^{\mu_t}\binom{\mu_t}{j}Z^j
       \sum_{k\ne0}\binom{-k-1}{j-1}Y_k\right).
\]
The global product of signs is retained.  This polynomial is still
homogeneous of degree $\ell$ in $Y$, and the marker $Z$ records
the weight of each homogeneous state component.  Apply
\eqref{eq:algebraic-rectangular-differences} while storing a
univariate array of degree at most $h$.  For
$a^{\langle i\rangle}$, apply after coefficient extraction the
linear functional
\[
 Z^w\longmapsto\binom{h+w}{i}.
\]
For $a$ itself, sum the marker coefficients.  This avoids expanding
$T(p_\mu)$ into all of its monomials.  It also handles the
inhomogeneous transform without confusing state weight with the
mode shift $r$.

All generalized binomial coefficients have arguments of magnitude
$O(W)$ and bit length $O(W)$.  Products contain at most $W$
factors, and the factorials and falling factorials above have
polynomial bit length.  The univariate arrays have polynomial
rank.  Thus every evaluation, accumulation, and exact division in
this part uses polynomial space and polynomial-bit integers.

\paragraph{Streaming the Schur basis changes.}
For partitions $\lambda,\mu$ of weight $w\le W$, put
$e_j=m_\mu(j)$ and set power sums outside
$\operatorname{supp}(\mu)$ to zero.  Then
\[
 [p_\mu]s_\lambda
 =\frac1{\prod_j e_j!}
   \sum_{0\le z_j\le e_j}
   (-1)^{\sum_j(e_j-z_j)}
   \prod_j\binom{e_j}{z_j}\,s_\lambda(p_j=z_j).
\]
Here weighted homogeneity suffices: a different monomial of
weight $w$ cannot be componentwise at least the target, because
any strict increase would have weight greater than $w$.
Therefore it has a below-target coordinate and vanishes under
that finite difference.  The grid has
$\prod_j(e_j+1)=2^{O(\sqrt W)}$ points.

At each grid point compute the complete functions by
\[
 m h_m=\sum_{j=1}^m p_jh_{m-j},\qquad h_0=1,
\]
and evaluate the Jacobi--Trudi determinant for $s_\lambda$ using
polynomial-time rational arithmetic.  Factorial common
denominators bound the bit lengths polynomially.  Hall
orthogonality gives the reverse coefficient
\[
 [s_\lambda]p_\mu=z_\mu[p_\mu]s_\lambda,
 \qquad z_\mu=\prod_jj^{m_\mu(j)}m_\mu(j)!.
\]

To compute a Schur mode entry, stream over the power-sum
partitions for its state, input, and output, multiply the three
basis-conversion coefficients by the power-sum mode entry just
computed, and accumulate.  For a transformed Schur state use
linearity, $T(s_\lambda)=\sum_\mu[p_\mu]s_\lambda\,T(p_\mu)$,
and the marker procedure.  Each of the three partition sets has
size $2^{O(\sqrt W)}$, so this fixed number of nested enumerations
and coefficient computations still takes $2^{O(\sqrt W)}$ time.
The rational basis coefficients have common denominators dividing
weight factorials; during their finite-difference computation a
fixed power of $W!$ is a sufficient common denominator.  Thus
accumulating exponentially many terms does not create
exponential bit lengths.

Finally, obtain a factorization entry from its indicated mode
entry and sign.  To obtain an entry of a full crossing piece,
stream its polynomially many $i$ values and the intermediate
partitions.  This is only a further fixed number of the same
enumerations.  Theorem~\ref{thm:integral-crossing} proves that
the final coordinates are integers.  Any final exact division
is performed over the integers before specialization to a target
ring.  No large table is materialized: partition indices, the
derivative vector, one finite-difference point, polynomial-size
univariate arrays, and accumulators suffice.  This proves all
three stated bounds.
\end{proof}

\section{Ordered insertion}
% Copyright 2026 icekylinx and GPT-6 Astra Pro. Apache-2.0.
% Curated from sat_research_checkpoint.zip: paper/direct_mixed_crossing_transfer.tex
% Complete simultaneous-insertion and ordered-crossing lemmas retained; downstream field/older numerical branches omitted.
\subsection{A direct noncommutative transfer in every characteristic}
\label{sec:direct-mixed-crossing}

The full crossing identity contains more information than its
commutative sum-of-squares contraction. Keeping that identity permits a
direct transfer with no division by two. In particular, the integral
Schur construction also gives noncommutative circuit bounds in
characteristic two.

\begin{lemma}[Simultaneous insertion without division]
\label{lem:simultaneous-insertion}
Let $R$ be a commutative unital ring. Suppose a division-free
commutative circuit of size $C$ simultaneously computes bilinear forms
$F_i(X,Y)$, and a division-free noncommutative circuit of size $D$
computes polynomials $G_i$. If
$F_i=\sum_{a,b}f_{i,a,b}X_aY_b$, then
\[
 \sum_i\sum_{a,b}f_{i,a,b}X_aG_iY_b
\]
has a division-free noncommutative circuit of size $O(C+D+t)$, where
$t$ is the number of forms. The statement is valid over $R$, including
in characteristic two and in the presence of zero divisors.
\end{lemma}

\begin{proof}
Keep the four bidegree components $(0,0),(1,0),(0,1),(1,1)$ of the
first circuit; this costs a constant factor and uses no division.
Its bilinear multiplication gates can be written
$A_\ell(X)B_\ell(Y)$, where both factors are linear. Regard each such
product as a new input $W_\ell$. The remaining bilinear computation is
a scalar-linear circuit
\[
 F_i=\sum_\ell c_{i,\ell}W_\ell.
\]
Transpose only this scalar-linear circuit to compute
$H_\ell=\sum_i c_{i,\ell}Z_i$. Explicitly, initialize the adjoints of
the output gates by $Z_i$, propagate an adjoint to both inputs of an
addition gate, and multiply by the same scalar when traversing a scalar
multiplication gate backwards. Sum the contributions arriving at each
gate. The number of operations is linear in the original number of
edges and the output count, hence $O(C+t)$.

Substitute $Z_i=G_i$ and output
$\sum_\ell A_\ell(X)H_\ell B_\ell(Y)$. Since every $c_{i,\ell}$
is central, this is the desired ordered insertion. The linear forms
and the transposed circuit remain shared. All identities use only
addition and multiplication, so the argument is characteristic-free.
\end{proof}

\begin{lemma}[Mixed crossing to an ordered identity polynomial]
\label{lem:direct-mixed-crossing}
Suppose arrays over $R$ satisfy
\begin{equation}
\label{eq:direct-mixed-crossing}
 [x=v][y=u]
 =\sum_{j\in J} f_{j,x,y}g_{j,v,u}
  +\sum_{j\in K} h_{j,x,u}k_{j,v,y}.
\end{equation}
If all associated bilinear forms are simultaneously computable by
division-free commutative circuits of total size $C$, then
\[
 \operatorname{ID}_{p,q}(X,Y)
   =\sum_{x\in[p],\,y\in[q]}X_xY_yX_xY_y
\]
has a division-free noncommutative circuit of size
$O(C+|J|+|K|)$.
\end{lemma}

\begin{proof}
Contract~\eqref{eq:direct-mixed-crossing} against the ordered word
$X_xY_yX_vY_u$. A row-oriented term is the product
\[
 \left(\sum_{x,y}f_{j,x,y}X_xY_y\right)
 \left(\sum_{v,u}g_{j,v,u}X_vY_u\right).
\]
A column-oriented term is
\[
 \sum_{x,u}h_{j,x,u}X_x
   \left(\sum_{v,y}k_{j,v,y}Y_yX_v\right)Y_u.
\]
Bihomogenization turns a commutative bilinear circuit into either
ordered bilinear circuit with constant overhead. When one shared
factor circuit has both row-oriented $XY$ outputs and column-oriented
$YX$ outputs, make two ordered copies and select the required outputs
from each. This costs a constant factor; no shared product is flipped
independently for different outputs. All column insertions
together are handled by Lemma~\ref{lem:simultaneous-insertion}; row
products and their sums add only the number of terms. The right side
therefore contracts to $\operatorname{ID}_{p,q}$ with the stated cost.
No polarization, square root of $-1$, or inversion of two is used.
\end{proof}

\section{The full three-shape integral matrix theorem}
% Copyright 2026 icekylinx and GPT-6 Astra Pro. Apache-2.0.
% Curated from sat_research_checkpoint.zip: paper/integral_full_rectangular_transfer.tex
% Complete full-U mathematical deduction retained with U expanded locally; historical formal-build claims and all weaker numerical/fallback corollaries omitted.
\subsection{The full rectangular curve from integral character profiles}
\label{sec:integral-full-rectangular}

We combine a recent integral coefficient interface with the weighted character-profile analysis of a separate source.  Zhou and
Wang~\cite{zhouwang2026} construct an integer tensor semiring,
prove an integral convolution bound and shared-leg separation,
and derive a universal circuit theorem for the symmetric
rectangular curve.  Their public Lean endpoint covers that
symmetric curve.  The following deduction applies the weighted
profile argument of Amanov, van den Berg, and
Zuiddam~\cite[Sections~4.5--4.6]{amanov2026} to the integral
ingredients, before the optional pairing of characters in
Zhou and Wang's numerical analysis. Their mathematical note also gives a weaker, unformalized bound for
three arbitrary sorted dimensions: $A+B$ when $B\ge2C$, and
$A+B+(2C-B)^2/(4(A+C-B))$ otherwise. The deduction here
strengthens that existing universal three-shape bound to the full
curve $U$; it is not a claim to be the first universal-ring treatment
of arbitrary rectangular shapes.

\begin{theorem}[Additional-source universal integer rectangular curve]
\label{thm:integral-full-rectangular}
Retain the integral spectrum, convolution, separation, and
sector constructions of~\cite{zhouwang2026}.  For every fixed
$A\ge B\ge C\ge0$ and $\epsilon>0$, there is one family of
ordered integer bilinear circuits of size
\[
 O\bigl(n^{U(A,B,C)+\epsilon}\bigr)
\]
for the corresponding rectangular matrix product, valid in
every associative unital ring.  The family and its bound are
chosen before the ring.  The function is explicitly
\[
 U(A,B,C)=\begin{cases}
 A+B,&C(A+B)\le AB,\\[2pt]
 \dfrac{(AB+BC+CA)^2}{4ABC},&C(A+B)>AB,
 \end{cases}
\]
where the first branch includes $C=0$, including $A=B=C=0$.  Integer scalar operations can be
implemented using $0,1,-1$ with the same slackened exponent.
\end{theorem}

\begin{proof}
We spell out the exact interface needed from the integral
source.  Let $\mathcal S_{\mathbb Z}$ contain zero and the
integer three-tensors that restrict to the unit tensor by
integer linear maps, modulo mutual restriction.  Direct sum
and tensor product give its semiring operations.  Its integer
restriction preorder satisfies the Strassen hypotheses, and
the source proves the duality
\[
 \widetilde R_{\mathbb Z}(T)
   =\sup_{\lambda\in\Delta(\mathcal S_{\mathbb Z})}\lambda(T)
 \qquad(T\in\mathcal S_{\mathbb Z}).
\]
All tensors used below have a unit restriction.  In particular,
polynomial convolution, oriented dot products, the sector sums,
and the coefficient-extraction tensors lie in this semiring.
There is no assertion that every nonzero integer tensor lies
in it: a scalar multiple such as $2xyz$ need not do so.

For a character $\lambda$, let $p_i\in[0,1]$ be the exponent
of the oriented dot product whose singleton leg is $i$.
The usual integer reindexings give
$\lambda(B_i(q))=q^{p_i}$ and factor the matrix tensor into
its three oriented dot products.  Write
\[
 C(r,s)=\sum_{0\le i<r,\,0\le j<s}x_i y_j z_{i+j}.
\]
The integral source supplies the following statements for
each individual character, without first averaging it with a
permuted character:
\begin{enumerate}
\item $1\le\lambda(C(r,s))\le r+s-1$, and the same for
every permutation of its legs.  The upper bound comes from
the integral asymptotic-rank convolution lemma, using
Vandermonde adjugates, algebra norms, and a finite B\'ezout
combination after tensor powering.  It is not interpolation
with divisions in the target ring.
\item Character values are monotone under integer coordinate
restrictions and under the displayed integral degenerations.
For the latter, coefficient extraction after tensor powering
has only polynomial rank overhead.
\item If a matched-sector tensor $T$ shares leg $i$ and has
sectors $T_j$, then
\[
 \log\lambda(T)\ge p_i H(q)+\sum_jq_j\log\lambda(T_j)
\]
for each probability vector $q$.  For $p_i>0$, optimization
over $q$ gives
\[
 \lambda(T)^{1/p_i}\ge\sum_j\lambda(T_j)^{1/p_i}.
\]
\item The unimodular determinant filtration degenerates
$C(r,s)\otimes B_1(2)$ into the matched sectors
$C(r,s+1)$ and $C(r,s-1)$, for $s\ge2$.
The factor $B_1(2)$ has value $2^{p_1}$, so the preceding
inequality yields exactly
\[
 \lambda(C(r,s+1))^{1/p_1}
 +\lambda(C(r,s-1))^{1/p_1}
 \le2\lambda(C(r,s))^{1/p_1}.
\]
The factor two is therefore the required discrete-concavity
constant, not a loss that needs to be removed.
\item The integer three-sector construction degenerates
$C(r,3s+r-1)$ into two copies of $C(r,s)$ and one copy with
the last two legs exchanged, all sharing the first leg.
Its middle-branch first-index reversal is an integer
permutation.  Applying the preceding entropy inequality gives
the two coupled tripling inequalities.  Every leg permutation
has the same properties.
\end{enumerate}

These are precisely the character-profile hypotheses used in
Lemma~4.10 and Sections~4.5--4.6 of~\cite{amanov2026}.
To make their field independence explicit, fix a first
coordinate and two corresponding positive row profiles
$g,h$, raised to the power $1/p_i$ when $p_i>0$.
They are discretely concave and satisfy
\[
 2g(s)+h(s)\le g(3s+r-1),\qquad
 2h(s)+g(s)\le h(3s+r-1).
\]
The weighted increment inequality of
\cite[Lemma~4.13]{amanov2026} is a real-variable consequence
of these two facts.  For positive integers $u,v,d,e$ it says
\[
 \left(\frac{g(u+d)}{g(u)}\right)^e
 \left(\frac{h(v+e)}{h(v)}\right)^d
 \ge
 \left(\frac{D}{D-4de/(d+e)}\right)^{d+e},
 \qquad D=u+v+d+e+r-1.
\]
It follows by bounding both increments below by a common
future secant slope and applying weighted AM--GM; no operation
on the coefficient ring occurs.

For positive integers $x_1,x_2,x_3$, put
$S=x_1+x_2+x_3$ and
$Q=p_1x_2x_3+p_2x_1x_3+p_3x_1x_2$.
Apply the weighted inequality to the six permuted convolution
profiles on the integer ray $(nx_1,nx_2,nx_3)$, pairing the
permutations that fix a first leg.  With
$\lambda^\pi(T)=\lambda(\pi\cdot T)$, their weighted product is
\[
 G_n=\prod_{\pi\in S_3}
 \lambda^\pi\!\left(C(nx_{\pi(1)},nx_{\pi(2)})\right)^{x_{\pi(3)}}.
\]
It obeys
\[
 \left(\frac{n+1}{2}\right)^{8Q/S}
       \le G_n\le(nS)^{2S}.
\]
The lower bound is the telescoping product of the six
increment bounds; the upper bound uses only the first
interface item.  When a parameter $p_i$ is zero, monotonicity
supplies the corresponding factor one, exactly as in
\cite[Proposition~4.14]{amanov2026}.
Letting $n$ grow, then using homogeneity and continuity, proves
\begin{equation}\label{eq:integral-character-quadratic}
 4(p_1x_2x_3+p_2x_1x_3+p_3x_1x_2)
       \le(x_1+x_2+x_3)^2
 \quad\text{for all }x_i\ge0.
\end{equation}

The weighted matrix character exponent is
$Ap_2+Bp_3+Cp_1$ under a consistent leg labeling.
If $C(A+B)\le AB$, substitute
$(x_1,x_2,x_3)=(A+B,B,A)$ into
\eqref{eq:integral-character-quadratic}, and use $p_1\ge0$;
the result is $Ap_2+Bp_3+Cp_1\le A+B$.
In the other branch, substitute $(AB,BC,CA)$ to obtain
\[
 Ap_2+Bp_3+Cp_1
       \le\frac{(AB+BC+CA)^2}{4ABC}=U(A,B,C).
\]
Thus every integer character satisfies the full three-shape
bound.  Integral spectral duality supplies the same
asymptotic integer-rank bound.

Finally fix positive slack before selecting one finite integer
rank witness.  For rational shape exponents choose a common
base so that all three dimensions are integer powers, then
take a sufficiently large fixed tensor block.  Recursing on
this ordered bilinear integer scheme costs
$O((t+1)\rho^t)$ at depth $t$: the rank $\rho$ is at least
each of the three input/output coordinate dimensions by
flattening over $\mathbb Q$, and all scalar-linear operations
are charged.  Integer coefficients in this fixed scheme are
implemented by addition and negation.  Padding and rational
approximation from above cover real exponents; continuity of
$U$ and the reserved slack pay for these steps and the
logarithmic recursion factor.  The resulting integer circuit
is chosen before its coefficient ring and preserves the
left-before-right order of each multiplication.
\end{proof}

% Copyright 2026 Yilin Liu and GPT-6 Astra Pro. Apache-2.0.
% Curated connective exposition of the checkpoint's rectangular extraction,
% direct mixed-crossing transfer, and integral full-curve numerical consequence.
\section{Integral rectangular extraction and shared evaluation}
\label{sec:rectangular-hurwitz}
For integer dimensions $p,q\ge2$, choose the least $h,E$ with
$p(h)\ge p$ and $p(E)\ge q$, where $p(t)$ denotes the partition
function in this paragraph only. Use any $p$ partitions of $h$ and $q$
partitions of $E$ in Theorem~\ref{thm:integral-crossing}. Restriction to
these coordinates preserves its integer crossing identity. For each
piece select its smaller rank factorization. The number of resulting
rank-one terms is at most
\[
 (2(E+h)+1)(2h+1)p(E+h).
\]
The partition estimate
$\log p(t)=\pi\sqrt{2t/3}+O(\log(t+2))$
\cite{hardyramanujan1918} therefore gives
\[
 P\le\exp\!\left(\sqrt{(\log p)^2+(\log q)^2}\right)
      (\log(pq+2))^{O(1)}.
\]
Here logarithms are natural. The rank-one terms have the two orientations
of \eqref{eq:direct-mixed-crossing}; the selected factorizations are
integer, with normalization one. No conversion to squares or division
by two is made. Their entries have polynomial logarithmic bit length
and are explicitly computable by
Lemma~\ref{lem:algebraic-coefficient-computation}.

Fix $a,b>0$ and $L\ge\max\{a,b\}$, and set
$s=\sqrt{a^2+b^2}$. Choose base dimensions
$p=n^{a+o(1)}$, $q=n^{b+o(1)}$. Then
$P=n^{s+o(1)}$. Partition each target variable list into
$U=n^{L-a+o(1)}$ and $V=n^{L-b+o(1)}$ blocks, respectively.
Replicate the same integer crossing decomposition for every block pair,
using the same pair at both repeated occurrences. This yields an exact
mixed crossing for $ID_{n^{L+o(1)}}$; padding adds zero coordinates.

For a factor family with coefficient matrices $A_j\in\mathbb Z^{p\times q}$,
all forms
\[
 (X^{(r)})^{\mathsf T}A_jY^{(t)},
 \quad j\in[P],\ r\in[U],\ t\in[V],
\]
are computed by two shared matrix products. Stack the $A_j$ into an
$(Pp)\times q$ constant matrix and multiply by the $q\times V$ matrix
of $Y$ blocks. Reshape the result to $p\times(PV)$, then multiply on
the left by the $U\times p$ matrix of $X$ blocks. Both factor families
require only a constant number of such calls. If their integer matrix
schemes have full arithmetic exponents $W_1,W_2$, the two shapes are
\[
 (s+a,b,L-b),\qquad (L-a,a,s+L-b),
\]
and Lemma~\ref{lem:direct-mixed-crossing} gives exponent
$\max\{W_1,W_2\}/L$ for the ordered identity. All $PUV$ output forms
are charged by the matrix schemes' output dimension, so insertion's
linear overhead is absorbed. Matrix products are used first as
commutative bilinear circuits in formal variables; ordered insertion
then preserves the required word order. The integer scheme, including
all scalar-linear operations, is fixed before evaluating any target
associative ring.

When $a+b=L$, the shapes are permutations of
$(s+a,b,a)$. Tensor-leg permutations of an integer bilinear matrix
identity give valid ordered schemes: interpret the first chosen leg
as the left input and the second as the right input, without commuting
target-ring values. This is the specialization used below.

\section{The selected constant-free ordered circuit bound}
\begin{theorem}[One integer circuit family with exponent $1.719516$]
\label{thm:selected-crossing}
Using the integral source constructions and weighted-profile theorem
specified in Theorem~\ref{thm:integral-full-rectangular}, there is one
family of noncommutative arithmetic circuits over $\mathbb Z$ computing
\[
 ID_N=\sum_{i,j=1}^{N}X_iY_jX_iY_j
\]
with $O(N^{1.719516})$ gates and only scalar constants $0,1,-1$.
The same chosen circuit evaluates correctly in every associative
unital ring, including characteristic two and rings with zero divisors.
The bound counts arithmetic gates; it does not bound the bit cost of
arbitrary ring-element evaluations or give a circuit-generation time
matching this gate exponent.
\end{theorem}
\begin{proof}
Set
\[
 a=\frac{167}{400},\qquad b=\frac{233}{400},\qquad L=1,
 \qquad s=\frac{\sqrt{82178}}{400}.
\]
Both shared matrix shapes are permutations of $(s+a,b,a)$. Put
\[
 R=\frac{286667054264}{10^9},\qquad A=a+R/400.
\]
Integer squaring gives $R^2>82178$. Padding the largest dimension
therefore permits the rational shape $(A,b,a)$, with $A>b>a>0$.
It lies in the curved region $a(A+b)>Ab$, and exact rational
substitution gives
\[
 \frac{1719516}{10^6}
 -\frac{(Ab+ba+aA)^2}{4Aba}
 =\frac{358038466018273}{1379112402223945625000}>0.
\]
Theorem~\ref{thm:integral-full-rectangular} supplies one integer
ordered matrix scheme at this rational shape, with any fixed
positive exponent slack. Reserve part of the strict displayed gap
for that slack. The shared-products construction and
Lemma~\ref{lem:direct-mixed-crossing} supply the ordered $ID_N$
circuit with the same exponent, apart from the partition factors,
rounding and padding. Their polynomial-logarithmic overhead fits
another part of the same gap.

There are two kinds of scalar constants. Crossing entries use
$h,E=O((\log(N+2))^2)$ and have bit length
$(\log(N+2))^{O(1)}$ by
Lemma~\ref{lem:algebraic-coefficient-computation}. For the matrix
scheme, fix the finite integer witness after fixing the positive
slack. Its finite collection of integer constants is independent of
$N$. Recursive evaluation, bihomogenization, scalar-linear
transposition and insertion use only these constants and $0,1,-1$.
If a scalar-only subcircuit appears, retain it instead of replacing
it by a new unanalyzed scalar leaf.

An integer of bit length $d$ is built by $O(d+1)$ doublings, additions
and a possible negation from $0,1,-1$. Replacing every scalar leaf
therefore costs only a polynomial-logarithmic factor. Reserve the
remaining strict gap for this cost. All constructions are identities
in the free associative algebra over $\mathbb Z$. Under evaluation in
any associative unital ring integer scalars map to central elements,
so the exact same circuit remains valid. No target-ring inverse or
field-extension constant is used.
\end{proof}

\paragraph{Inherited output reductions.}
The degree-four distinct-index polynomial
$S_{4,N}=\sum_{i_1,i_2,i_3,i_4\text{ distinct}}X_{i_1}X_{i_2}X_{i_3}X_{i_4}$
and the row-ordered rectangular polynomial
$\operatorname{perm}_{4,N}=\sum_{i_1,i_2,i_3,i_4\text{ distinct}}
 X_{1,i_1}X_{2,i_2}X_{3,i_3}X_{4,i_4}$
have the same bound by the established integer identities of
Hrube\v{s}~\cite[Proposition~26]{hrubes2026}. These are output transfers
of the same construction, not separate algorithmic contributions.
For completeness put $p_j=\sum_i X_i^j$, $ID^*=ID_N(X,X)$, and
\[
 S_2=p_1^2-p_2,\qquad
 S_3=p_1S_2-p_2p_1-\sum_i X_i p_1 X_i+2p_3.
\]
Then $S_4=p_1S_3-A_0-B_0-C_0$, where
\begin{align*}
 A_0&=p_2S_2-p_3p_1-\sum_iX_i^2p_1X_i+2p_4,\\
 B_0&=\sum_i X_i p_1X_i p_1-ID^*
       -\sum_iX_i p_1X_i^2-p_3p_1+2p_4,\\
 C_0&=\sum_iX_iS_2X_i-\sum_iX_i^2p_1X_i
       -\sum_iX_i p_1X_i^2+2p_4.
\end{align*}
Every sum outside $ID^*$ has linear circuit size, with integer
coefficients. Finally substitute for each $X_j$ the strictly
upper-shift $5\times5$ matrix whose $(r,r+1)$ entry is $X_{r,j}$.
The $(1,5)$ entry is exactly $\operatorname{perm}_{4,N}$.
Matrix simulation costs a constant factor per gate and needs no
division. This concerns four rows; it is not the ordinary growing-order
square permanent.

\bibliographystyle{plainnat}
\bibliography{references}
\end{document}
